\section{Introduction}

\subsection{Context and Motivation}

Cloud computing has fundamentally transformed the economics and operational models of data persistence, shifting infrastructure costs from fixed capital expenditure to variable operational expenditure metered at sub-second granularity \citep{haubenschild2025oltp, vanrenen2023cloud}. Within this landscape, serverless computing platforms—exemplified by AWS Lambda, Azure Functions, and Google Cloud Functions—have emerged as a dominant paradigm, abstracting server provisioning and enabling event-driven, elastically scalable application architectures \citep{bodner2025serverless, liu2023faaslight}. These serverless workloads impose stringent requirements on data-layer persistence: sub-millisecond latency, elastic throughput, and predictable per-operation costs without over-provisioned capacity.

Amazon DynamoDB, a fully managed NoSQL key-value and document database service, has become a canonical choice for serverless data persistence due to its integration with AWS Lambda, automatic horizontal scaling, and advertised single-digit millisecond latency at any scale \citep{elhemali2022dynamodb, idziorek2023transactions}. However, achieving these performance characteristics in practice depends critically on two configuration dimensions that remain underexplored in controlled experimental studies: \textbf{partition-key design} and \textbf{capacity mode}.

Partition-key design determines how data is distributed across DynamoDB's underlying storage partitions. A poorly chosen partition key can concentrate read and write traffic on a subset of partitions (a ``hot key'' or ``hot partition'' problem), leading to throttling even when table-level capacity is sufficient \citep{aws2025adaptive}. Three design patterns are documented in vendor guidance and grey literature: simple keys (high-cardinality unique identifiers such as order IDs), composite keys (partition key plus sort key, enabling range queries and customer-localised access patterns), and write-sharded keys (partition key with appended shard suffix to distribute write-heavy traffic) \citep{aws2025sharding, eldeeb2022neuroshard}. Yet empirical performance and cost comparisons across these patterns under realistic serverless workloads remain absent from peer-reviewed literature.

Capacity mode governs billing and admission control. DynamoDB offers two modes: \textbf{on-demand} (per-request pricing with automatic scaling) and \textbf{provisioned} (reserved RCU/WCU with optional auto-scaling and lower unit costs) \citep{aws2025capacitymode}. The on-demand mode simplifies capacity planning at the expense of higher per-request costs, while provisioned mode requires accurate throughput forecasting but offers cost savings for predictable workloads. The interaction between partition-key design and capacity mode—for example, whether write sharding remains beneficial under on-demand billing, or whether provisioned tables with poor key design experience greater throttling—has not been systematically studied.

Existing benchmarking literature, including the YCSB suite and cloud analytics benchmarks \citep{ferreira2025ycsb, vanrenen2023cloud}, typically evaluates NoSQL databases in self-hosted or fixed-capacity scenarios. \citet{pantelic2026benchmarking} provide a recent, methodologically rigorous baseline, benchmarking SQL and NoSQL persistence (including DynamoDB-equivalent workloads) under read-heavy, write-heavy, and mixed microservice profiles on a single self-hosted node. Their study quantifies latency and throughput but, by design, omits metered cost, throttling, and consumed capacity units—variables that only exist in cloud-metered infrastructure. Similarly, studies of cloud-native databases such as Aurora Serverless \citep{barnhart2024aurora} and disaggregated storage systems \citep{pang2024disaggregated, durner2023object} focus on architectural principles rather than configuration-level decision heuristics for practitioners.

This gap motivates the present research: a controlled empirical evaluation that simultaneously varies partition-key design and capacity mode on Amazon DynamoDB under Lambda-driven serverless workloads, measuring latency, throttling, consumed RCU/WCU, and cost per 10,000 operations.

\subsection{Research Question and Objectives}

\textbf{Research Question:} How does Amazon DynamoDB table configuration—specifically the interaction between partition-key design and capacity mode—determine the performance–cost trade-off under serverless workloads?

This study pursues four specific objectives:

\begin{enumerate}
    \item \textbf{Objective 1:} Quantify how partition-key design (simple, composite, write-sharded) affects latency and throttling rate under each workload profile (read-heavy, write-heavy, mixed, burst).
    \item \textbf{Objective 2:} Quantify how capacity mode (on-demand vs. provisioned with auto-scaling) affects throttling rate and cost at equivalent load.
    \item \textbf{Objective 3:} Determine whether the latency-minimising configuration also minimises cost, or whether trade-offs exist.
    \item \textbf{Objective 4:} Express each of the six configurations (3 key designs × 2 capacity modes) as a position on the trade-off surface spanned by latency, throttling rate, and cost per 10,000 operations.
\end{enumerate}

\subsection{Contribution}

This study makes three principal contributions to the cloud database systems literature:

\begin{enumerate}
    \item \textbf{Experimental framework for configuration-level trade-offs:} A controlled factorial design and reproducible artefact (Terraform, Lambda drivers, Zipfian workloads, CloudWatch-oriented metrics schema, ANOVA pipeline) targeting partition-key design $\times$ capacity mode on Amazon DynamoDB. Live metered campaign results are \textbf{not} yet collected; evaluation evidence in this report is moto simulation / placeholders only, and must not be read as empirical findings on production DynamoDB.
    \item \textbf{Reusable artefact:} A complete Infrastructure-as-Code artefact (Terraform, Lambda handlers, Zipfian workload generators, CloudWatch metrics collectors, statistical analysis scripts) is provided, enabling third-party replication and extension to other DynamoDB configurations or workload distributions.
    \item \textbf{Practical guidance for cloud architects:} Findings are expressed as actionable heuristics for serverless application developers: under which workload characteristics (read/write mix, access skew, burst vs. sustained load) does each configuration minimise latency, cost, or a weighted combination thereof.
\end{enumerate}

\subsection{Limitations and Scope}

This study evaluates Amazon DynamoDB only; findings may not generalise to other NoSQL services (Azure Cosmos DB, Google Cloud Firestore) due to differences in partition algorithms, billing granularity, and auto-scaling behaviour. The experiment uses synthetic order/customer data (1,000,000 items, 10,000 customers) with 1 KB items and Zipfian access distributions; real-world workloads with different item sizes, data models, or access patterns may exhibit different trade-offs. Finally, this study measures steady-state and burst behaviour over minutes to hours; diurnal or seasonal traffic patterns at weekly or monthly timescales are out of scope.

\subsection{Report Structure}

The remainder of this report is organised as follows. Section 2 reviews related work on cloud database benchmarking, partition-key design, capacity management, and serverless persistence. Section 3 describes the experimental design, including the factorial structure (3 × 2 × 4), workload profiles, dependent variables, and statistical analysis plan. Section 4 details the artefact implementation: Terraform infrastructure, Lambda workload generators, metrics collection, and cost modelling. Section 5 presents experimental results, including latency distributions, throttling rates, consumed capacity, and cost analysis. Section 6 discusses findings, interprets trade-offs, and provides configuration selection heuristics. Section 7 concludes with contributions, limitations, and future work.
